\documentclass[11pt,a4paper]{article}
\usepackage[margin=1in]{geometry}
\usepackage{hyperref}
\usepackage{enumitem}
\usepackage{xcolor}
\usepackage{listings}

\lstset{
  basicstyle=\ttfamily\small,
  breaklines=true,
  frame=single,
  columns=fullflexible,
}

\title{\bfseries Week 1: Project Scaffolding, Frontend
  Setup \& Vercel CI/CD}
\author{Yash \\ \small Roll No: 1024030846}
\date{27 August 2026}

\begin{document}
\maketitle

\section{Context}
Start of the SkillGrid project. The GitHub repository
existed but was empty. Today's work was to bring the
repository up to the structure required by the UCS503P
course template, and to get a first, real, automatically
deployed increment of the frontend working -- in line
with the course's emphasis on rapid, CI/CD-backed
iterations.

\section{Repository Structure}
Cloned the reference course template
(\texttt{tiet-ucs503/ucs503p-202627odd-template}) to
study its layout, then replicated the same structure in
the SkillGrid repository:

\begin{itemize}[leftmargin=0.7cm]
  \item \texttt{docs/} -- MkDocs-based documentation site.
  \item \texttt{journals/} -- one folder per team member
    (this file lives under \texttt{journals/1024030846-yash/}).
  \item \texttt{project-proposal/},
    \texttt{project-report-prototype-stage/},
    \texttt{project-report-final/} -- LaTeX report folders.
  \item \texttt{code/} -- source code root.
  \item \texttt{assets/}, \texttt{mkdocs.yml},
    \texttt{.github/workflows/mkdocs.yml} -- documentation
    build and GitHub Pages deployment.
\end{itemize}

\subsection{Issue: Symlinks on Windows}
The template links \texttt{docs/assets} and
\texttt{docs/journals} to the root-level
\texttt{assets/} and \texttt{journals/} folders via
symlinks, so MkDocs can serve them without duplicating
content.

\textbf{Problem:} creating these with \texttt{ln -s} on
Windows (without Developer Mode/admin privilege) silently
created real directories instead of symlinks, which then
got fully duplicated into the git index as separate
binary files.

\textbf{Fix:} removed the duplicated directories and used
git's plumbing commands directly to record proper symlink
blob entries (git file mode \texttt{120000}):

\begin{lstlisting}[language=bash]
printf '../assets' | git hash-object -w --stdin
git update-index --add --cacheinfo 120000,<sha>,docs/assets
\end{lstlisting}

This checks out as an inert placeholder file on Windows
(no symlink support), but resolves as a real symlink on
the Ubuntu-based GitHub Actions runner, which is where the
MkDocs build actually happens.

\section{Frontend Scaffold}
Since \texttt{code/} was empty, scaffolded a minimal
Next.js application (TypeScript, ESLint, Tailwind CSS,
App Router) under \texttt{code/frontend/}, so that the
deployment pipeline would have something real to build and
verify end-to-end, rather than wiring up CI/CD against
nothing.

Verified locally before wiring up any automation:
\begin{lstlisting}[language=bash]
npm install
npm run build
\end{lstlisting}
Both succeeded, confirming the scaffold was a valid
starting point.

\section{CI/CD Pipeline: Vercel Deployment}
Added \texttt{.github/workflows/vercel-deploy.yml}: a
GitHub Actions workflow that deploys \texttt{code/frontend}
to Vercel using the official Vercel CLI flow
(\texttt{vercel pull} $\to$ \texttt{vercel build} $\to$
\texttt{vercel deploy --prebuilt}), triggered on pushes to
\texttt{main} (production) and on pull requests (preview),
scoped to only run when \texttt{code/frontend/**} changes.

\subsection{Secrets}
The workflow needs three GitHub Actions secrets:
\texttt{VERCEL\_TOKEN}, \texttt{VERCEL\_ORG\_ID}, and
\texttt{VERCEL\_PROJECT\_ID}. Noted that Vercel tokens are
scoped to an account/team, not to an individual project --
the actual per-project targeting comes from
\texttt{VERCEL\_PROJECT\_ID}, so the same token can be
reused for a backend project added later.

Confirmed via \texttt{gh secret list} that all three secrets
were present on the repository before the first CI run.

\subsection{Debugging: Duplicated Root Directory Path}
\textbf{Error encountered on first real deploy run:}
\begin{lstlisting}
Error: ENOENT: no such file or directory, open
'.../code/frontend/code/frontend/package.json'
\end{lstlisting}

\textbf{Relevant context:} the Vercel project's own
\emph{Root Directory} setting (configured when the project
was imported on vercel.com) was \texttt{code/frontend}.
The workflow \emph{also} set
\texttt{defaults.run.working-directory: code/frontend}
before invoking the Vercel CLI.

\textbf{Key observation:} \texttt{vercel build} applies the
pulled project's Root Directory setting \emph{relative to
the current working directory}. Since the workflow had
already \texttt{cd}'d into \texttt{code/frontend}, Vercel
appended the same subpath again, producing
\texttt{code/frontend/code/frontend}.

\textbf{Solution:} removed the
\texttt{working-directory: code/frontend} default from the
workflow, so all Vercel CLI steps run from the repository
root and the project's own Root Directory setting is the
single source of truth:
\begin{lstlisting}[language=diff]
- defaults:
-   run:
-     working-directory: code/frontend
\end{lstlisting}

\emph{Because} the two mechanisms (shell working directory
and CLI-applied project setting) are not aware of each
other and are not idempotent when composed -- only one of
them should determine the subdirectory.

\section{Outcome}
After the fix, the workflow ran green end-to-end (install
$\to$ pull $\to$ build $\to$ deploy) and the frontend went
live at:
\begin{center}
  \url{https://skillgrid-frontend-coral.vercel.app}
\end{center}

\section{Follow-ups}
\begin{itemize}[leftmargin=0.7cm]
  \item \texttt{mkdocs} GitHub Pages workflow also failed
    on its first run and needs a separate look.
  \item \texttt{project-proposal/main.tex} for SkillGrid
    itself is still a TODO.
  \item Backend architecture and its own Vercel/CI setup
    not yet decided.
\end{itemize}

\end{document}
